\chapter{Allgemeine Sicherheit}\label{ch:gefahrdungsbeurteilung}

\begin{table}[H]
    \centering
    \renewcommand{\arraystretch}{1.5}

    \begin{tabularx}{\textwidth}{
        |>{\centering\arraybackslash}m{0.27\textwidth}
        |>{\centering\arraybackslash}X
        |>{\centering\arraybackslash}m{0.27\textwidth}|}
        \hline

        \begin{minipage}[c][3.0cm][c]{\linewidth}
            \centering
            \textbf{Bereich:}\\[0.6em]
            Gefahren die aus dem Modell hervorgehen
        \end{minipage}
        &
        \begin{minipage}[c][3.0cm][c]{\linewidth}
            \centering
            \Large\textbf{Gefährdungsbeurteilung}\\
            in Anlehnung an §§ 5,6 ArbSchG
        \end{minipage}
        &
        \begin{minipage}[c][3.0cm][c]{\linewidth}
            \centering
            \textbf{Projekt:}\\[0.6em]
            Zweiachsig nachgeführtes PV-Modell
        \end{minipage}
        \\

        \hline
    \end{tabularx}\label{tab:table}
\end{table}

\begin{GPMTable}
{L{3.2cm}|L{5.0cm}|L{2.2cm}| Y}
{}
{}
{\TableHead{Gefährdung} &
\TableHead{Gefahr} &
\TableHead{persöhnliche Beurteilung vor Maßnahme} &
\TableHead{Gegenmaßnahmen}}


    \multicolumn{4}{|l|}{\textbf{Mechanische Gefährdungen}} \\
    \hline

    &
    Quetschstellen an beweglichen Teilen &
    Akzeptabel &
    Bewegliche Teile vor der Arbeit abschalten.\\
    \hline

    &
    Verletzungen durch scharfe Kanten &
    Inakzeptabel &
    Bauteile auf scharfe Kanten prüfen und diese gegebenenfalls entschärfen.\\
    \hline

    &
    Einzugsgefahr an Motorwellen &
    Inakzeptabel &
    Motoren vor Arbeit von Spannung trennen. \\
    \hline

    &
    Kippen oder Umstürzen der Anlage &
    Akzeptabel &
    Von Tischkanten fernhalten und auf waagerechtem Untergrund stellen. \\
    \hline


    \multicolumn{4}{|l|}{\textbf{Elektrische Gefährdungen}} \\
    \hline

    &
    Berührung spannungsführender Teile &
    Inakzeptabel &
    Anlage vor Arbeiten spannungsfrei schalten und Spannungsfreiheit prüfen. \\
    \hline

    &
    Kurzschluss nach oder bei Arbeiten an der Verdrahtung &
    Inakzeptabel &
    Verdrahtung vor der Inbetriebnahme kontrollieren und geeignete Absicherung verwenden. \\
    \hline

    &
    Beschädigte Leitungen oder Isolierungen &
    Inakzeptabel &
    Regelmäßige Sichtprüfung und Austausch beschädigter Leitungen. \\
    \hline

    &
    Gefahr durch Kondensatoren &
    Inakzeptabel &
    Vor arbeiten Kondensatoren entladen. \\
    \hline

    &
    Überlastungsgefahr &
    Inakzeptabel &
    Sicherungen einbauen. \\
    \hline


    \multicolumn{4}{|l|}{\textbf{Chemische Gefährdungen}} \\
    \hline

    &
    Hautkontakt mit gesundheitsschädlichen Stoffen &
    Inakzeptabel &
    Nach Berührung der Elemente Händewaschen. \\
    \hline

    &
    Einatmung von gesundheitsschädlichen Dämpfe &
    Inakzeptabel &
    Nur in belüfteten Räumen oder mit Abzug arbeiten. \\
    \hline

    \multicolumn{4}{|l|}{\textbf{Brand und Explosionsgefahr}} \\
    \hline

    &
    Zündquellen durch Lichtbögen &
    Inakzeptabel &
    Alle Kontakte vor einschalten Prüfen. \\
    \hline

    &
    Erwärmung von Bauteilen &
    Akzeptabel &
    Sicherungen zur Verhinderung von Überlastung, sowie Bauteile auf die Belastung auslegen. \\
    \hline


    &
    Explosion von Kondensatoren &
    Inakzeptabel &
    Kondensatoren auf korrekte Polung prüfen. \\
    \hline


\end{GPMTable}
\newpage
\begin{UnifiedInfoBox}
    \textbf{Allgemeine Maßnahmen und der Umgang bei Arbeiten an der Anlage}

    \begin{itemize}
        \item Arbeiten an der elektrischen Anlage grundsätzlich nur im spannungsfreien Zustand durchführen.
        \item Die Anlage vor Beginn der Arbeiten gegen unbeabsichtigtes Einschalten sichern.
        \item Bewegliche Teile und mögliche Quetschstellen während der Arbeiten beachten.
        \item Nur zertifizierte und geeignete Werkzeuge verwenden und beschädigte Werkzeuge nicht einsetzen.
        \item Leitungen, Steckverbindungen und Bauteile vor der Inbetriebnahme auf Beschädigungen prüfen.
        \item Den Arbeitsbereich sauber und frei von unnötigen Gegenständen halten.
        \item Änderungen an der Verdrahtung oder Steuerung vor dem Einschalten kontrollieren.
        \item Die Anlage bei ungewöhnlichem Verhalten, Geruch, Geräuschen oder Erwärmung sofort abschalten.
        \item Praktische Arbeiten nur mit mindestens zwei Personen durchführen.
        \item Kein offenes Schuhwerk tragen.
        \item Keine Nahrung in der Arbeitsumgebung aufbewahren.
        \item Nach Abschluss der Arbeiten alles sicher verstauen und zurückräumen.
    \end{itemize}

\end{UnifiedInfoBox}
